%% kapa: exam=1 round=37 year=2026 period=1 form=A subject=real_estate subject_ko=부동산학원론 seq=32 number=112
%% source: Q-Net 감정평가사 기출문제 (cst00309) artlSeq=5258924 file=raw/1008/5258924/2026년 제37회 감정평가사 1차 1교시.pdf
%% conversion: kordoc 4.24.51 -> tools/gen1.py (auto); review=proofread-sonnet; flags=none
\begin{question}{112}
민간투자 사업방식에 관한 설명으로 옳은 것을 모두 고른 것은?
\begin{framedtext}
ㄱ. BTO방식은 기반시설의 완공과 동시에 국가나 지방자치단체인 공공에 소유권이 귀속되는 방식으로, 운영과 관리에 대한 책임이 공공에 있다. \par
ㄴ. BOT방식은 민간사업자가 투자비를 조달하여 기반시설을 건설한 후 일정기간 운영하고 사업기간이 종료할 때 공공에 소유권을 이전하는 방식으로, 주로 대규모 프로젝트에서 사용된다. \par
ㄷ. BLT방식은 민간사업자가 투자비를 조달하고 기반시설을 건설한 후 운영권을 임대하여 투자비를 회수하고, 약정기간 만료 후에 소유권이 공공에 이전된다. \par
ㄹ. BOO방식은 민간사업자가 기반시설을 건설한 후 공공에 소유권을 이전하는 방식으로, 투자비 회수가 어려워 사전 목표 수익률이 확보되어야 한다. \par
\end{framedtext}
\begin{choices}
  \item ㄱ, ㄹ
  \item ㄴ, ㄷ
  \item ㄱ, ㄴ, ㄹ
  \item ㄴ, ㄷ, ㄹ
  \item ㄱ, ㄴ, ㄷ, ㄹ
\end{choices}
\end{question}
\answer{2}
